\section{Les batteries lithium-soufre}

\setcounter{equation}{0}
\setcounter{Q}{0}

%Mines-Pont, Chimie 2019 - Filière PSI
La technologie des batteries rechargeables est en constante évolution pour répondre à des exigences de plus en plus fortes en termes de densité d'énergie, vitesse de recharge, limitation de l'auto-décharge, durabilité et sécurité. Ces exigences parfois contradictoires poussent à leur limite certains matériaux qui subissent des transformations importantes lors des cycles, sans oublier l'échauffement important qui peut se produire. Parmi les multiples systèmes prometteurs étudiés par la communauté scientifique, les batteries lithium-soufre ont attiré une attention particulière ces dernières années en raison d'une densité d'énergie théorique beaucoup plus élevée que celle des batteries Li-ion classiques et d'un coût plus faible, le soufre étant bon marché.

\Q Etablir la configuration électronique de l'atome de soufre dans l'état fondamental. Quel anion peut-il former ?

La structure cristallographique de la phase stable du soufre à température et pression ambiantes est orthorhombique, la maille élémentaire est un parallélépipède rectangle de paramètres $a=1.0 ~ \si{\nano\meter}$, $b=1.3 ~ \si{\nano\meter}$ et $c=2.4 ~ \si{\nano\meter}$. La maille élémentaire contient $16$ unités \ce{S_k}. La masse volumique de cette phase est d'environ $2.1 ~ \si{\gram\per\centi\meter\cubed}$.

\Q Déterminer k.

De façon schématique, les batteries lithium-soufre sont constituées d'une anode de lithium, d'une cathode de soufre élémentaire (dispersé dans du carbone conducteur) et d'un électrolyte organique permettant le transfert des ions \ce{Li^+} jusqu'à la cathode où ils précipitent avec les sulfures résultant de la réduction du soufre.

\Q Ecrire l'équation de la réaction de décharge et calculer la tension standard en circuit ouvert de la batterie à \SI{298}{\kelvin}. Commenter sachant que la tension nominale de batteries Li-ion est de l'ordre de $3.6 ~ \si{\volt}$.

\Q Calculer la capacité spécifique théorique (en \si{\ampere\hour\per\gram}) du matériau de cathode. Comparer avec des matériaux typiques de cathode de batteries Li-ion dont la capacité spécifique théorique est de l'ordre de \SI{300}{\milli\ampere\hour\per\gram}.

\Q En supposant le lithium et le soufre en proportions st{\oe}chiométriques, évaluer la variation relative du volume total des phases solides entre l'état totalement chargé et l'état totalement déchargé. Commenter brièvement ce résultat.

La courbe de la première décharge de la batterie a l'allure suivante :

\Q Quelle fraction de la capacité spécifique théorique récupère-t-on ?

En réalité, le soufre ne se réduit pas directement sous forme de sulfure simple mais de polysulfures \ce{S_n^{2-}}, avec $2 \leq n \leq 8$. Ces espèces sont solubles dans l'électrolyte. Dans le cadre d'un modèle simplifié, on considère que le premier plateau (noté 1) correspond à la réduction de l'espèce \ce{S_8^0} en solution (résultant de la dissolution du soufre solide) en un seul type de polysulfure \ce{S_p^{2-}} (en réalité il s'agit d'une réduction progressive via des polysulfures longs). Le deuxième plateau (noté 2) correspond à la réduction de \ce{S_p^{2-}} en \ce{S^{2-}} et \ce{S_2^{2-}}.

\Q Déterminer $p$ et en déduire l'écriture des réactions correspondant à chaque plateau.

Dans l'électrolyte considéré, les longueurs d'onde $\lambda_{max}$ du maximum d'absorption de lumière UV-visible des espèces du soufre en solution sont données dans le tableau suivant.

\begin{tabular}{cccccccccc}
$\lambda_{max} ~ [\si{\nano\meter}]$ & 255 & 265 & 280 & 332 & 344 & 425 & 435 & 470 & 490 \\
espèce & \ce{S^{2-}} & \ce{S_8^0} & \ce{S_2^{2-}} & \ce{S_6^{2-}}  & \ce{S_3^{2-}} & \ce{S_4^{2-}} & \ce{S_5^{2-}} & \ce{S_7^{2-}} & \ce{S_8^{2-}}
\end{tabular}

\Q Proposer un protocole pour doser les espèces du soufre en solution tout au long de la décharge. On en précisera les limitations.

\Q Proposer une structure de Lewis pour le polysulfure linéaire \ce{S_4^{2-}}.

La mesure de l'absorbance à \SI{425}{\nano\meter} en fonction du taux de décharge donne le résultat suivant :

\begin{center}
\includegraphics[width=10cm]{Ex08/Fig02}
\end{center}

\Q Discuter la cohérence de ce résultat avec le modèle simplifié utilisé pour expliquer les plateaux de décharge.
